\section{Discussion}
\label{sec:implications}

\paragraph{Readability belongs to the probe, ownership to the reader.}
Checkpoints sharing a vision tower have identical probes and write vectors, yet own different concepts, and the chest-trained towers read and answer findings best in the grid but own only \cfSpecTowerOwned{} of \cfSpecTowerN{} chest cells. Ownership therefore belongs on the scorecard of the checkpoint to be deployed rather than on that of the vision tower it contains.

\paragraph{On chest radiographs, the label direction and the answer direction come apart.}
On natural images, the label direction and the model's answer direction are substantially aligned, and writing the label direction controls the answer. On chest radiographs the two are nearly orthogonal, and the answer direction is owned in \cfAnsChestOwnedA{} of \cfAnsChestN{} chest cells where the label direction is owned in \cfAnsChestOwnedLabel{}. The flat rows of the chest write matrices fit a reader that responds to a signal shared across findings rather than to each finding's own direction \citep{pahde2025navigating}; which signal that is remains open. On chest radiographs a legible label does not mark the direction along which the reader chooses its answer.

\paragraph{What a steering evaluation should report.}
Probe accuracy and answer accuracy, the two quantities most often reported for medical VLMs, do not predict whether a write is specific to the concept it names. Random and sham writes show that a direction reaches the answer; only directions fitted the same way for the competing findings, written at the same site and steering strength in the same patients, show whether the direction belongs to its label. Studies that use probe normals as interventions should report the write matrix $W$, compare $O_q$ across layers, strengths, and models, and pair every clinical write with a natural-image control, where the same procedure does produce ownership.

